\section{带电粒子在磁场中的运动}

\subsection{洛伦兹力}
在本章初我们曾经提到过，阴极射线管中的电子束在磁场力的作用下将发生偏转，本节将讨论单个点电荷在运动时所受的磁场力，这种运动电荷在磁场中受到的力称为\uwave{洛伦兹力}
\begin{BoxEquation}[洛伦兹力]
    \vb*{F}=q\vb*{v}\times\vb*{B}
\end{BoxEquation}
\Refx{公式}{洛伦兹力}是实验结论，其表明，运动电荷受到的磁场力的方向总是垂直于其速度和磁感强度构成的平面，具体的指向还与运动电荷自身的电荷量的正负有关，即
\begin{itemize}
    \item 对于正电荷，其磁场力是右手四指由速度转向磁感强度时拇指指向的正方向。
    \item 对于负电荷，其磁场力是右手四指由速度转向磁感强度时拇指指向的负方向。
\end{itemize}
由于洛伦兹力总是垂直于电荷的运动方向，因此，\textbf{洛伦兹力只能改变运动电荷的运动方向}，而并不改变运动电荷的速率和动能，这也相当于是说洛伦兹力不对运动电荷做功。

\subsection{带电粒子在匀强磁场中的运动}
现在我们来研究一种特殊情况，即带电粒子在匀强磁场中的运动。为了便于分析，我们将带电粒子的速度分解为平行于磁感强度和垂直于磁感强度两个分量，随后再将运动合成。
\begin{Figure}{带电粒子在匀强磁场中的运动}
    \inputfigure[scale=1]{Chapter10D_Figure01}
\end{Figure}
如果我们记速度和磁感强度的夹角为$\theta$，平行分量即为$v\sin\theta$，垂直分量即为$v\cos\theta$。

对于平行于磁感强度的速度分量$\vb*{v}\cos\theta$，其并不会导致任何的洛伦兹力，因此在平行于磁感强度的方向上，带电粒子将保持进入磁场前的$v\cos\theta$的速度做匀速直线运动。

对于垂直于磁感强度的速度分量$\vb*{v}\sin\theta$，其会使得带电粒子受到一个垂直于该速度分量的洛伦兹力，而且由于这里的磁场是匀强的，其受到的洛伦兹力的大小并不改变，因此在垂直于磁感强度的方向上，带电粒子将以$v\sin\theta$的速度做匀速圆周运动、如果我们将这两个方向的运动叠加合成，我们会看到，\textbf{带电粒子在匀强磁场中的运动轨迹为一条等速螺旋线}。

\begin{Figure}{带电粒子在匀强磁场中的运动分解}
    \subfigure[主视图]
    {
        \inputfigure[scale=1]{Chapter10D_Figure03}
    }\qquad\qquad
    \subfigure[俯视图]
    {
        \inputfigure[scale=1]{Chapter10D_Figure02}
    }
\end{Figure}

现在我们来具体分析一下这个运动，设带电粒子的质量为$m$，设带电粒子的电荷量为$q$，假设圆周运动的半径为$R$，根据\Refx{公式}{洛伦兹力}和\Refx{公式}{向心加速度}
\begin{Equation}
    F=qv\sin\theta B=\frac{mv^2\sin^2\theta}{R}
\end{Equation}
这样就解得了匀速圆周运动的半径，称为\uwave{回旋半径}
\begin{BoxEquation}[回旋半径]
    R=\frac{mv\sin\theta}{qB}
\end{BoxEquation}
进而我们就可以解得周期
\begin{equation*}
    T=\frac{2\pi R}{v\sin\theta}=\frac{2\pi mv\sin\theta}{qBv\sin\theta}=\frac{2\pi m}{qB}
\end{equation*}
这样就算得了匀速圆周运动的周期，称为\uwave{回旋周期}，其倒数称为\uwave{回旋频率}。
\begin{BoxEquation}[回旋周期]
    T=\frac{2\pi m}{qB}
\end{BoxEquation}
\begin{BoxEquation}[回旋频率]
    \nu=\frac{qB}{2\pi m}
\end{BoxEquation}
\Refx{公式}{回旋周期}指出，带电粒子的回旋周期与其速度无关，其仅取决于磁感强度和带电粒子的质量与电荷量，也就是说，速度的大小和方向只会改变回旋半径和稍后将介绍的回旋螺距，不影响带电粒子的回旋周期，这一结论是磁聚焦和回旋加速器的理论基础。

\Refx{公式}{回旋周期}中我们看到，回旋周期与磁感强度和电荷量成正比，回旋周期与质量成反比，这一结果可以这样理解，较强的磁场意味着电荷更容易被磁场“拽住”，从而回旋的空间就小一些了,另外，更大的电荷量意味着磁场力更强，更大的质量意味着运动状态更难被磁场力改变。有时也会将$m/q$称为带电粒子的\uwave{质荷比}，这样就可以简单的说，\textbf{回旋周期与质荷比和磁感强度成正比}，不过理论中更为常用的是$q/m$即带电粒子的\uwave{荷质比}。

带电粒子在磁场中回旋一周就会前进一定距离，这一距离称为\uwave{回旋螺距}，不难想象，回旋螺距的大小即为带电粒子在平行磁场方向的速度$v\cos\theta$与回旋周期$T$的乘积，即
\begin{BoxEquation}[回旋螺距]
    H=2\pi\frac{mv\cos\theta}{qB}
\end{BoxEquation}
\Refx{公式}{回旋螺距}和\Refx{公式}{回旋半径}中，我们看到，\textbf{回旋半径和回旋螺距均与质荷比和磁感强度成正比}，但两者同时也与速度成正比，并且和速度的方向有关
\begin{itemize}
    \item 速度一定时，速度与磁感强度的夹角越大，回旋半径越大。
    \item 速度一定时，速度与磁感强度的夹角越小，回旋螺距越大。
\end{itemize}
由此可见，某种意义上，速度的方向决定了速度在回旋半径和回旋螺距上的分配。

这里有一个简单的应用，我们知道，虽然我们将电子束假象为一条射线，但实际上电子的角度总归会存在一定的微小差异，这就会使得电子束在光屏上的影像变得模糊，这时我们就可以沿着电子束的方向施加一个匀强磁场，设电子束中的电子的运动速度均为$v$，但与电子束的整体方向存在一个夹角，设电子的运动方向与磁场的夹角为$\theta$，由于这里的$\theta$很小
\begin{equation*}
    v\cos\theta=v\qquad v\sin\theta=v\theta
\end{equation*}
这里我们看到，角度不同的电子在垂直磁场方向的分量$v\theta$不同，但是这些电子平行磁场的分量$v$基本相同，而各个电子的回旋周期又相同，这就意味着，尽管各个电子在磁场中以不同的回旋半径进行运动，但是，在时间上相隔一个回旋周期，在空间上相隔一个回旋螺距，这些电子会重新会聚到同一位置，这样就可以显著的改善电子束的成像，由于这种现象与光束经过透镜后聚焦的现象类似，因此将这种通过磁场会聚带电粒子的方法称为\uwave{磁聚焦}。

下面我们来看一些通过带电粒子在匀强磁场中的运动特性制造的实验设备。

\subsection{回旋加速器}
\uwave{回旋加速器}是原子核物理学中获得高速粒子的一种装置，这种装置的基本原理就是利用回旋周期与速率无关的性质，回旋加速器的核心部分是两个D形盒，它的形状类似于扁圆的金属盒沿直径剖开得到的左右两半，在D形盒之间留有狭缝并存在交变的电场，在D形盒内部由于导体空腔对电场的屏蔽效应，几乎不存在电场。整个装置被置于一个巨大的真空容器中，在垂直D形盒的方向上用电磁铁产生一个均匀磁场，覆盖整个D形盒空间。

当带电粒子最初位于狭缝时，通过电场使其取得一个速度$v_1$并进入右侧的D形盒，在磁场的洛伦兹力作用下以回旋半径$R_1$经半个圆周回到狭缝，这时期望电场恰好反向，从而可以继续加速带电粒子，使其以更大的速度$v_2$在左侧D形盒中以$R_2$为半径再次回旋。
\begin{Figure}{回旋加速器的最初两轮回旋}
    \subfigure[回旋加速器的第一轮加速]
    {
        \inputfigure[scale=0.45]{Chapter10D_Figure04}
    }\qquad
    \subfigure[回旋加速器的第二轮加速]
    {
        \inputfigure[scale=0.45]{Chapter10D_Figure05}
    }
\end{Figure}
根据\Refx{公式}{回旋半径}，我们看到
\begin{equation*}
    R_1=\frac{mv_1}{qB}\qquad R_2=\frac{mv_2}{qB}
\end{equation*}
回旋加速器在设计上的巧妙之处在于，虽然回旋半径随着速度的增大而不断增大，但是回旋周期却并不变化，这也就是说，为了保证带电粒子在每一次经过狭缝时都得到加速，只需要使用一个交流频率为带电粒子的回旋频率的电源就可以了，而不需要人为动态的根据粒子的速度调节电场方向。在交流电场的作用下，带电粒子会在狭缝中被不断加速并取得越来越大的回旋半径，带电粒子最终会接近D形盒的边缘，将其引出后就可以进行相关实验。
\begin{Figure}{回旋加速器}
    \inputfigure[scale=0.7]{Chapter10D_Figure06}
\end{Figure}
回旋加速器能使得带电粒子取得的最大速度，取决于其D形盒半径的大小，因此如果我们需要获得速度很高的带电粒子，就需要建造庞大的回旋加速度器，其半径尺度可能达到千米量级，另外，质荷比较小的带电粒子可以在半径一定的回旋加速器中取得更大的速度。

回旋加速器也存在一定缺陷，这就是当带电粒子的速度接近光速时，由于相对论效应的作用，带电粒子的质量开始增大，带电粒子的质荷比也开始增大，这就会改变带电粒子的回旋频率，这时交流电源的频率与之不相匹配，从而无法进一步加速带电粒子（因为这时带电粒子经过狭缝区域时，有时电场方向与运动方向相同，有时电场方向与运动方向相反）。

如果要进行更进一步的加速，就需要使用原理不同的\uwave{同步加速器}。世界上最大的同步加速器是大型强子对撞机\footnote{简称LHC，即Large Hadron Collider。}，其位于一个圆周达到27公里的环形隧道内，贯穿法国和瑞士边境。

\subsection{速度选择器}
\uwave{速度选择器}是一种可以控制通过的带电粒子速度的设备，它由两块带电平行金属板产生的均匀电场$\vb*{E}$和与之垂直的均匀磁场$\vb*{B}$构成，从粒子源进入速度选择器的带电粒子将会同时受到电场力$\vb*{F}_e=qE$和磁场力$\vb*{F}_m=qvB$的作用，两个力的方向相反，因此
\begin{itemize}
    \item 如果电场力大于磁场力，那么带电粒子将坠向右侧的负极板，无法通过速度选择器。
    \item 如果磁场力大于电场力，那么带电粒子将坠向左侧的正极板，无法通过速度选择器。
\end{itemize}
由此可见，那些可以通过速度选择器的带电粒子必然受到相同的电场力和磁场力。
\begin{Figure}{速度选择器}
    \inputfigure[scale=0.7]{Chapter10D_Figure07}\vspace{-5pt}
\end{Figure}\vspace{-5pt}
由于电场力和磁场力应相等，故
\begin{equation*}
    qE=qvB
\end{equation*}
这就是说，只有速度满足电场和磁场之比的带电粒子才能无偏转的通过速度选择器
\begin{BoxEquation}[速度选择器]
    v=\frac{E}{B}
\end{BoxEquation}
这样我们就可以通过控制电场和磁场的强度之比，来获取我们所需的特定速度的带电粒子。

\subsection{质谱仪}
同种元素的质子数是一定的，但可以具有不同的中子数，因此可以有不同的原子量，这就是所谓的\uwave{同位素}，例如对于氢元素，它具有三种同位素：氕（\ce{_1^1H}）、氘（\ce{_1^2H}）、氚（\ce{_1^3H}）。
\begin{Figure}{氢元素的同位素}
    \subfigure[氕（\ce{_1^1H}）]
    {
        \inputfigure[scale=2]{Chapter10D_Figure08}
    }\qquad
    \subfigure[氘（\ce{_1^2H}）]
    {
        \inputfigure[scale=2]{Chapter10D_Figure09}
    }\qquad
    \subfigure[氚（\ce{_1^3H}）]
    {
        \inputfigure[scale=2]{Chapter10D_Figure10}
    }\qquad
\end{Figure}
同位素的测定可以通过\uwave{质谱仪}实现，首先我们要将待测的同位素混合物电离，使其以离子状态存在，随后这些离子需要经过速度选择器，假设控制的通过速度为$v$，根据\Refx{公式}{回旋半径}，由于经过速度选择器的离子具有相同的速度，因此离子在磁场中的回旋半径就只与离子的质荷比有关了，质荷比越大回旋半径越大，质荷比越小回旋半径越小。
\begin{Figure}{质谱仪}
    \inputfigure[scale=0.7]{Chapter10D_Figure11}
\end{Figure}
如果我们在离子进入磁场回旋半周的位置放置一块感光胶片（或类似传感器），由于不同质荷比的同位素具有不同的回旋半径，因此同位素离子落在胶片上的位置就不同。根据胶片的感光位置就可以推演出各个同位素的质量值，根据感光强弱还可以推出同位素的比例。

如果记质谱仪中的磁场为$\vb*{B}$，而记速度选择器中的磁场和电场为$\vb*{B}_0,\vb*{E_0}$，根据\Refx{公式}{回旋半径}和\Refx{公式}{速度选择器}，经过质谱仪的离子的回旋半径满足以下公式
\begin{BoxEquation}[质谱仪]
    R=\frac{m}{qB}\frac{E_0}{B_0}
\end{BoxEquation}

\subsection{霍尔效应}
实验表明，当我们将一块矩形导体板垂直放置在匀强磁场中时，在矩形的长度方向上通以电流时，在矩形的宽度方向的两个面间会存在电势差，这种现象称为\uwave{霍尔效应}。
\begin{Figure}{霍尔效应}
    \subfigure[正载流子的霍尔效应]
    {
        \inputfigure[scale=0.65]{Chapter10D_Figure12}
    }\qquad
    \subfigure[负载流子的霍尔效应]
    {
        \inputfigure[scale=0.65]{Chapter10D_Figure13}
    }
\end{Figure}
现在我们来研究霍尔效应是如何产生的，我们先考虑载流子为正电荷的情况，正电荷原先依照电流方向在横向做定向运动，但是在磁场$\vb*{B}$的作用下，正电荷会在纵向上发生漂移，这样的结果便是，在导体的1面将富集正电荷，在导体的2面将对应的带负电荷，从而在导体内部形成了一个附加电场$\vb*{E}_H$，称为\uwave{霍尔电场}，其导致的电势差就称为\uwave{霍尔电势差}。

由于霍尔电场力$F_e=qE_H$与磁场力$F_m=qvB$的方向相反，因此霍尔电场的出现会减弱正电荷在纵向上的漂移，随着霍尔电场的增强，最终霍尔电场力和磁场力将达到平衡
\begin{equation*}
    qE_H=qvB
\end{equation*}
这时正电荷也不再发生纵向漂移，因此在导体中就形成了一个稳定的霍尔电场
\begin{BoxEquation}[霍尔电场的速度表示]
    E_H=vB
\end{BoxEquation}
这时的霍尔电场是匀强的，如果记1,2面间的距离为$l$，那么1,2面间的霍尔电势差即为
\begin{BoxEquation}[霍尔电势差的速度表示]
    U_H=vBl
\end{BoxEquation}
这样就解释了霍尔效应产生的原因，不过在通常情况下我们并不知道导体中载流子的运动速度，而我们知道的是导体中通过的电流，因此我们需要将以上两个公式用电流表示。

根据\Refx{公式}{电流的微观表述}
\begin{equation*}
    v=\frac{I}{nqS}
\end{equation*}

这里的$S$是导体的横截面积，我们已经设导体的宽度为$l$，如果记导体的厚度为$d$，那么
\begin{equation*}
    v=\frac{I}{nqld}
\end{equation*}

上式代入\Refx{公式}{霍尔电场的速度表示}和\Refx{公式}{霍尔电势差的速度表示}即得
\begin{Equation}[霍尔效应1]
    E_H=\frac{IB}{nqld}\qquad U_H=\frac{IB}{nqd}
\end{Equation}
上式中的$1/nq$是仅与导体材料有关的系数，因为$n$是载流子的数密度而$q$是载流子的电荷量，因此我们可以将$1/nq$定义为一个关于材料的系数，称为\uwave{霍尔系数}，记作$R_H$
\begin{BoxDefinition}[霍尔系数]
    R_H=\frac{1}{nq}
\end{BoxDefinition}
运用\Refx{定义}{霍尔系数}，式\ref{式:霍尔效应1}就可以表述为
\begin{BoxEquation}[霍尔电场]
    E_H=\frac{IB}{nqld}=\frac{IB}{ld}R_H
\end{BoxEquation}
\begin{BoxEquation}[霍尔电势差]
    U_H=\frac{IB}{nqd}=\frac{IB}{d}R_H
\end{BoxEquation}
由此我们可以看到，当材料的霍尔系数越大，当施加的电流和磁感强度越强，当使用的导体越薄，霍尔效应就越明显，霍尔电势差与导体宽度无关，霍尔电场则与导体宽度成反比。

以上讨论是对于载流子为正电荷的情况进行的，而对于载流子为负电荷的情况，这时电荷的正负和电荷的运动方向都反向了，因此电荷受到的洛伦兹力的方向是不变的，因而电荷仍然向着1面漂移，但由于这里为负电荷，此时在1面与2面间形成的电场会反向，然而电荷受到的电场力的方向仍然是不变的，因而最终电场力也会与磁场力达成平衡，形成霍尔效应。

这里我们发现了一个重要的现象，在过去，我们常常认为正电荷的定向运动与等量负电荷的反向运动是等效的，就电流确实如此，就霍尔效应中两者受到的磁场力和电场力也是如此，但是，两者在宏观上产生的霍尔电场和霍尔电势差却是刚好相反，也就是说
\begin{center}
    霍尔效应可以分辨导体中载流子的类型
\end{center}
换言之，正电荷的定向运动和负电荷的反向运动在电流上等效，但其霍尔效应是不同的。

这里我们可能会有疑问，我们知道导体中的载流子总是带负电的电子，那么为什么我们还要研究载流子带正电的霍尔效应？这是因为我们如果将研究范围从导体扩大到半导体时，我们会发现，不同类型的半导体材料将表现出不同的霍尔效应，这会有两种结果
\begin{itemize}
    \item 霍尔效应表现为负载流子的材料，称为\uwave{电子型半导体}，其载流子是带负电的电子。
    \item 霍尔效应表现为正载流子的材料，称为\uwave{空穴型半导体}，其载流子是带正电的\uwave{空穴}。
\end{itemize}
这里的空穴虽然在微观上仍然是电子的缺失形成的，但我们却必须认为空穴型半导体中是带正电的空穴在运动，而不能视作是带负电的电子在反向运动，因为只有这样才能满足霍尔效应的实验结果，因此空穴在这里并不是一个简单的等效物理模型，而是满足霍尔效应所必须的客观存在。所以可以说，\textbf{载流子的正负性完全是被霍尔效应的结果所定义的}。

\Refx{定义}{霍尔系数}中，正载流子的霍尔系数为正，负载流子的霍尔系数为负，若规定右手四指由磁场方向转向电流方向时，右手拇指指向的为2面，右手拇指背向的为1面，并定义霍尔场强和霍尔电势差的正值方向为由1面指向2面的方向，那么\Refx{公式}{霍尔电场}和\Refx{公式}{霍尔电势差}就可以普遍的适用正载流子和负载流子的情况了。

\Refx{定义}{霍尔系数}指出，霍尔系数与材料的载流子数密度是负相关的，导体中电子的浓度大霍尔效应其实并不不明显，相反，在半导体中载流子的浓度要低的多，在半导体中的霍尔效应也就要显著的多，因此大部分\uwave{霍尔器件}都是用半导体材料制成。

霍尔效应的应用是非常广泛的，除了判定材料的载流子类型，霍尔效应的另外一个重要用途就是可以测定磁场强度，如果已知电流$I$和材料的霍尔系数$R_H$以及材料厚度$d$，且测得霍尔电势差为$U_H$，磁场强度就可以根据\Refx{公式}{霍尔电势差}通过计算得出
\begin{equation*}
    B=\frac{U_Hd}{R_HIB}
\end{equation*}
磁传感器通常就是利用霍尔效应和霍尔器件制成的（这比测定电流受到的力要容易的多）。



